Firefly flashes demonstrate biology's hierarchical organisation from ecosystems to quantum mechanics. Species-specific flash patterns enable mate recognition and, in tropical swarms, synchronous displays visible across forests. Neural circuits generate these patterns by controlling oxygen flow through tracheal valves to specialised photocytes. Within these cells, luciferase catalyses luciferin oxidation with extreme efficiency, converting chemical energy to light with minimal heat. The photons themselves arise when excited electrons in oxyluciferin transition between quantum energy states, emitting at 560–590 nm. 
